\begin{center}\begin{tikzpicture}[font=\small,>=stealth,line join=round]
\draw[gray!40,rounded corners=3pt] (0,0) rectangle (14.2,8.3);\node at (7.1,7.700000000000001) {易错二：只经过(1,1)，还不能判定是幂函数};
\begin{scope}[shift={(1.8741935483870968,2.9045454545454543)},x=1.9354838709677418cm,y=0.8363636363636363cm]\draw[->,gray] (-0.4,0)--(2.7,0) node[right]{$x$};\draw[->,gray] (0,-1.5)--(0,4) node[above]{$y$};\draw[teal,thick,solid,domain=0:2.3,samples=180] plot(\x,{\x});\draw[orange!85!black,thick,dashed,domain=0:2.3,samples=180] plot(\x,{2*\x-1});\fill[orange!85!black] (1,1) circle(2.2pt);\end{scope}\node[anchor=west,align=left,text width=5.0cm] at (8.4,4.0) {青：$y=x$是幂函数\\[9pt]橙：$y=2x-1$不是\\[9pt]两者都过$(1,1)$\\[9pt]但橙线有常数项$-1$\\[9pt]必须核对整个函数式};
\node[text width=12.6cm,align=center] at (7.1,.65) {图中的反例说明：过公共点是必要条件，不能单独当作判定任意函数为幂函数的充分条件。};
\end{tikzpicture}\end{center}
